\begin{longtable}[c]{ |c|c|c|c|c| }
\caption{ \begin{math}Key\end{math} generated based on \begin{math}Key_{A}\end{math}}
\label{tab:extmemencr-key-possibility}
\\\hline 
\rowcolor{lightgray}
Block$_A$ & \begin{math}Key\end{math}  & \begin{math}Key\end{math} Length (bit) \\\hline
\endfirsthead
\hline
\rowcolor{lightgray}
Block$_A$ & \begin{math}Key\end{math}  & \begin{math}Key\end{math} Length (bit) \\\hline
\endhead
\endfoot
Yes       &   \begin{math}           Key_{A}\end{math}    &   256   \\\hline
No        &   \begin{math}           0^{256}\end{math}    &   256   \\\hline
\end{longtable}
\vspace{-1.5em}
{\bfseries Notes: }\\
``YES" indicates that the block exists; ``NO" indicates that the block does not exist; ``\begin{math}0^{256}\end{math}" indicates a bit string that consists of 256-bit zeros.
Note that using \begin{math}           0^{256}\end{math} as Key is not secure. We strongly recommend to configure a valid key.
